\chapter{Current Status and Preliminary Discussion}

\section{Purpose of this chapter}

Because part of the thesis work is still ongoing, this chapter intentionally
provides a \emph{status-oriented discussion} instead of a final conclusion
chapter. The goal is to make progress explicit, identify current evidence, and
highlight what remains to reach thesis completion.

\section{Status by work package}

Table~\ref{tab:status_matrix} summarizes the current maturity of each work
package.

\begin{table}[tbp]
\centering
\caption{Current status matrix (draft stage).}
\label{tab:status_matrix}
\begin{tabular}{@{}p{0.22\linewidth}p{0.2\linewidth}p{0.5\linewidth}@{}}
\toprule
Work package & Status & Evidence available in this draft \\
\midrule
WP1: Approximate GOOP solving (paper) & Completed & Formal method, implementation, and published-style quantitative results \cite{delasheras:2025}. \\
WP2: Duckietown integration (\texttt{patos}) & Functional prototype & End-to-end software architecture, manager/adapter/bridge loop, deployment scripts, and test workflow documentation. \\
WP3: Inverse estimation (\texttt{InverseGoopSolver}) & Exploratory prototype & Centralized/decentralized runners and candidate-based preference-order inference workflow with reproducible outputs. \\
\bottomrule
\end{tabular}
\end{table}

\section{What is already demonstrated}

At this stage, the thesis already supports the following claims with concrete
artifacts:

\begin{enumerate}
    \item Approximate lexicographic IBR over time can drastically reduce solve
    time versus tightly coupled baseline solves in representative GOOP scenarios.
    \item The same planning concepts can be integrated into a practical robotics
    stack with clear interfaces and runtime safety fallbacks.
    \item Preference-order estimation can be operationalized as a reproducible
    inverse workflow on top of the same planning backbone.
\end{enumerate}

These claims are currently strongest for algorithmic feasibility and software
architecture. They are weaker for broad statistical generalization, which
requires additional experiments.

\section{Main limitations of the current draft}

The current draft has four important limitations:

\begin{itemize}
    \item \textbf{Incomplete WP2 quantitative evaluation.} Runtime robustness and
    closed-loop quality metrics in broader multi-robot conditions are still being
    collected.
    \item \textbf{Incomplete WP3 statistical evaluation.} Preference-order
    estimation is implemented but not yet benchmarked at the scale needed for
    final claims.
    \item \textbf{Combinatorial growth in inverse search.} Full permutation
    testing does not scale to large objective sets.
    \item \textbf{Partial integration of inverse methods with embodied tests.}
    The full bridge from inverse estimation to real-robot execution remains
    ongoing.
\end{itemize}

\section{Preliminary interpretation}

Even with these limitations, the current evidence suggests that the thesis
pipeline is coherent:

\begin{itemize}
    \item the forward method is fast enough to support repeated replanning,
    \item the software integration makes this method executable in a realistic
    robotics loop,
    \item inverse estimation can reuse the same formal objects and code
    abstractions, reducing conceptual and implementation fragmentation.
\end{itemize}

This coherence is important: it indicates that final thesis value is likely to
come from \emph{integration quality} as much as from any single algorithmic
result.

\section{Risks to final completion}

The main risks between this draft and the final thesis submission are:

\begin{enumerate}
    \item delays in collecting stable WP2 multi-robot metrics,
    \item insufficient statistical depth for WP3 evaluation,
    \item scope creep in inverse-model variants before consolidating a primary
    benchmark protocol.
\end{enumerate}

Mitigation actions are listed in the roadmap chapter.

\section{Takeaway}

This draft already establishes a strong technical baseline and integrated
architecture. The final thesis effort is now concentrated on completion-quality
evidence: broader experiments, consolidated metrics, and final synthesis across
forward solving, deployment, and inverse estimation.
